% chapters/ch12_vnv_do178.tex — Chapter 12: V&V, DO-178B & Test Methodology
% Scope (PLAN.md): white/black/grey box, statement vs branch vs MC/DC coverage,
% DO-178B levels A-E, requirements traceability, ATP vs QTP, HIL testing,
% fault injection, defect lifecycle & root cause (5-why), AS9100.
% Budget 30 pp, mix 24 Q&A / 2 code / 4 concept.
% All snippets verified gcc -Wall -Wextra -std=c11 (zero warnings).

\chapter{V\&V, DO-178B \& Test Methodology}
\label{ch:vnv-do178}

\begin{partnote}
\textbf{Scope} --- white / black / grey box testing $\rightarrow$ coverage
(statement / branch / MC/DC) $\rightarrow$ DO-178B levels A--E $\rightarrow$
requirements traceability $\rightarrow$ ATP vs QTP $\rightarrow$ HIL testing
$\rightarrow$ fault injection $\rightarrow$ defect lifecycle \& root-cause
(5-why) $\rightarrow$ AS9100. This is the candidate's professional core ---
DO-178B Level A verification, 100\% requirements-to-test traceability, ATP/QTP
authoring, white/black-box firmware testing --- so it must be the most
fluent chapter of all.
\end{partnote}

Verification and Validation answers two questions: \emph{``did we build the
product right?''} (verification --- against the requirements) and \emph{``did we
build the right product?''} (validation --- it meets the real need). In
safety-critical avionics this isn't paperwork; it's the discipline that lets you
\emph{claim}, with evidence, that flight software won't kill anyone. Every
question here is about producing that evidence rigorously.

%=====================================================================
\section{Primer}
%=====================================================================

\subsection{White, black, and grey box}

\textbf{Black-box} testing exercises the software through its interface against
its \emph{requirements}, with no knowledge of the internal code --- it answers
``does it do what the requirement says?'' \textbf{White-box} (structural) testing
uses knowledge of the code structure to exercise paths, branches, and conditions
--- it answers ``is every part of the code actually tested?'' \textbf{Grey-box}
combines them: requirement-driven tests informed by some internal knowledge
(e.g.\ to target boundary cases). A rigorous process uses black-box tests derived
from requirements \emph{and} measures white-box coverage to prove the
requirements-based tests actually exercised all the code --- which is exactly
the DO-178B model.

\subsection{Coverage: statement, branch, MC/DC}

Structural coverage measures how thoroughly tests exercise the code, in
increasing strictness. \textbf{Statement coverage}: every line executed at least
once (weak --- misses untaken branches). \textbf{Branch (decision) coverage}:
every decision taken both true and false. \textbf{MC/DC} (Modified
Condition/Decision Coverage): for every \emph{condition} in a compound decision,
a test must show that condition \emph{independently} affecting the outcome ---
i.e.\ flipping just that condition flips the result, with the others held fixed.
MC/DC is required for the highest assurance level because it catches errors in
the boolean logic that branch coverage misses (a condition that's never actually
the deciding factor). It needs roughly $N+1$ tests for $N$ conditions, not the
$2^N$ of exhaustive testing.

\subsection{DO-178B levels and the V-model}

DO-178B assigns a \textbf{Design Assurance Level} by the consequence of failure:
\textbf{A} (catastrophic --- loss of aircraft) the most rigorous, through
\textbf{B} (hazardous), \textbf{C} (major), \textbf{D} (minor), to \textbf{E}
(no effect). Higher levels demand more objectives --- notably, MC/DC coverage is
required at Level A, decision coverage at B, statement at C. The process follows
the \textbf{V-model}: requirements $\rightarrow$ design $\rightarrow$ code down
the left arm, with verification (reviews, then unit/integration/system test) up
the right arm, each test level tracing back to the artefact it verifies. The
guiding principle is \textbf{traceability}: every requirement traces to design,
code, and tests, and every test traces back to a requirement --- nothing
untested, nothing untraceable.

\subsection{ATP/QTP, HIL, fault injection, defect lifecycle}

An \textbf{ATP} (Acceptance Test Procedure) is the formal, numbered, step-by-step
procedure --- each step with an expected result --- run to \emph{accept} a
product against its requirements; a \textbf{QTP} (Qualification Test Procedure)
qualifies it against the full spec/environment. \textbf{HIL} (Hardware-in-the-
Loop) testing runs the real flight computer against a \emph{simulated} aircraft/
environment in real time, so you can exercise scenarios (and failures) you can't
fly. \textbf{Fault injection} deliberately introduces faults (bad inputs,
corrupted data, hardware failures) to verify the system detects and handles them
safely. The \textbf{defect lifecycle} tracks a found defect from report
$\rightarrow$ triage $\rightarrow$ root-cause (e.g.\ \textbf{5-why}) $\rightarrow$
fix $\rightarrow$ re-test/regression $\rightarrow$ close, with traceability
throughout --- and \textbf{AS9100} is the aerospace quality-management standard
wrapping all of it.

%=====================================================================
\section{Q\&A Bank}
%=====================================================================

\tierbanner{Tier 1 --- Screening (phone / HR-technical)}%
{Two to four sentences.}

\question{Q1.1}{Verification vs validation?}
Verification asks ``did we build it right?'' --- does the software meet its
specified requirements (reviews, tests). Validation asks ``did we build the right
thing?'' --- do those requirements and the product meet the actual operational
need. V\&V together give confidence the system is both correct and fit for
purpose.

\question{Q1.2}{Black-box vs white-box testing?}
Black-box tests through the interface against requirements, ignoring internal
code; white-box uses knowledge of the code structure to exercise its paths,
branches, and conditions. Black-box checks behaviour; white-box checks that the
code is thoroughly exercised.

\question{Q1.3}{What is statement coverage?}
The fraction of executable statements (lines) run at least once by the tests.
It's the weakest structural metric --- 100\% statement coverage can still miss
untaken branches and faulty conditions.

\question{Q1.4}{What is branch (decision) coverage?}
Every decision point (if/while/switch) is exercised in both its true and false
outcomes. Stronger than statement coverage because it forces both directions of
every branch.

\question{Q1.5}{What is MC/DC and where is it required?}
Modified Condition/Decision Coverage requires showing that each condition in a
compound decision independently affects the outcome. It's required for DO-178B
Level A (the most safety-critical) software, because it catches boolean-logic
errors that branch coverage misses.

\question{Q1.6}{What do the DO-178B levels A--E mean?}
They're Design Assurance Levels set by failure severity: A = catastrophic
(most rigorous), B = hazardous, C = major, D = minor, E = no safety effect
(least). Higher levels require more verification objectives and stricter
coverage.

\question{Q1.7}{What is requirements traceability?}
A documented bidirectional link from each requirement to its design, code, and
test cases, and from each test back to a requirement. It proves every
requirement is implemented and tested, and that no code/test exists without a
reason --- the backbone of a certifiable process.

\question{Q1.8}{ATP vs QTP?}
An Acceptance Test Procedure is the formal step-by-step procedure (each step with
an expected result) to accept a product against its requirements; a Qualification
Test Procedure qualifies it against the full specification/environment.
Both are controlled documents producing recorded pass/fail evidence.

\question{Q1.9}{What is HIL testing?}
Hardware-in-the-Loop runs the actual embedded target connected to a real-time
simulation of the plant/aircraft, so the real firmware experiences realistic
(and fault) scenarios without flying. It bridges bench unit testing and flight
test.

\question{Q1.10}{What is fault injection?}
Deliberately introducing faults --- bad inputs, corrupted data, simulated
hardware failures --- to verify the system detects and responds to them safely.
It tests the error-handling and safety paths that normal testing rarely
exercises.

\tierbanner{Tier 2 --- Core technical (panel round)}%
{A short paragraph.}

\question{Q2.1}{Why isn't 100\% statement coverage enough for safety-critical
code?}
Because a statement can be covered without both outcomes of the decisions around
it being tested. \code{if (a) x = 1;} reaches 100\% statement coverage with a
single test where \code{a} is true --- but the \code{a == false} path (where
\code{x} keeps its old value) is never exercised, and a bug there is invisible.
Branch coverage forces both directions; MC/DC additionally forces each condition
to be shown to matter. Higher assurance demands stronger coverage precisely
because weaker metrics give false confidence.

\question{Q2.2}{Explain MC/DC with a concrete example.}
Take \code{if (a \&\& b)}. MC/DC requires, for each condition, a pair of tests
that differ \emph{only} in that condition and produce different outcomes. For
\code{a}: tests (a=1,b=1)$\rightarrow$true and (a=0,b=1)$\rightarrow$false show
\code{a} independently flips the result. For \code{b}: (a=1,b=1)$\rightarrow$true
and (a=1,b=0)$\rightarrow$false show \code{b} matters. Three tests --- \{11, 01,
10\} --- achieve MC/DC for two conditions ($N+1$), versus four for exhaustive.
The checker in the Coding Gym verifies exactly these independence pairs exist.

\question{Q2.3}{Walk the V-model and where each test level fits.}
Down the left arm: system requirements $\rightarrow$ high-level (software)
requirements $\rightarrow$ design $\rightarrow$ code. Up the right arm, each
level verifies its left-arm counterpart: \textbf{unit/low-level tests} verify the
code against the design/low-level requirements (with structural coverage);
\textbf{integration tests} verify components work together per the architecture;
\textbf{system/acceptance tests} verify the whole against the high-level/system
requirements (ATP). The horizontal links are traceability: every test traces to
the requirement it verifies, and gaps (a requirement with no test, code with no
requirement) are defects in the process itself.

\question{Q2.4}{How does requirements traceability actually work day to day?}
Each requirement gets a unique ID; design elements, code, and test cases
reference those IDs, and a \textbf{traceability matrix} (often in a tool ---
DOORS, or a managed spreadsheet/SAP QM record) maps requirements
$\leftrightarrow$ tests. You can then query: which tests cover requirement X (and
did they pass), and which requirements lack a test (a coverage gap). 100\%
requirements-to-test traceability means every requirement has at least one
passing test and every test justifies itself --- the claim on the resume, and the
thing an auditor checks first. The Coding Gym models the coverage query.

\question{Q2.5}{What's the difference between a requirements-based test and a
structural-coverage measurement, and why do you need both?}
A requirements-based (black-box) test checks the software \emph{does what it
should}; structural coverage (white-box) measures \emph{how much of the code} the
tests exercised. You need both because requirements tests can pass while leaving
code paths untested (dead code, unanticipated branches, error handlers), and
coverage without requirements just exercises code without checking correctness.
DO-178B's model is: derive tests from requirements, run them, then \emph{measure}
structural coverage --- and any uncovered code must be explained (it's dead code
to remove, or a missing requirement/test). The two together close the loop.

\question{Q2.6}{Why is HIL testing valuable, and what are its limits?}
HIL lets the real firmware experience realistic dynamic scenarios --- including
dangerous or rare failure conditions (engine-out, sensor failure, out-of-range
inputs) --- in real time, repeatably, without risking an aircraft, and far
earlier and cheaper than flight test. Its limits: the simulation is only as good
as its models (an unmodelled effect won't be caught), real-time fidelity and I/O
timing must be accurate, and it can't fully replace flight test for the
un-modelled real world. It's the crucial middle of the verification pyramid
between bench and flight.

\question{Q2.7}{Describe the defect lifecycle and the role of root-cause
analysis.}
A defect is \textbf{reported} (with steps to reproduce), \textbf{triaged}
(severity/priority, assigned), \textbf{root-caused} (find the actual cause, not
the symptom --- e.g.\ \textbf{5-why}: keep asking ``why?'' until you reach the
true origin), \textbf{fixed}, \textbf{re-tested} (the specific case) and
\textbf{regression-tested} (nothing else broke), then \textbf{closed} --- all
recorded with traceability to the requirement/test involved. Root-cause matters
because fixing the symptom leaves the real defect to recur; a 5-why on ``the
output was wrong'' might reach ``a shared variable wasn't protected'' (Chapter 4)
--- the actual fix.

\question{Q2.8}{What is AS9100 and how does it relate to DO-178B?}
AS9100 is the aerospace \emph{quality management system} standard (built on
ISO\,9001 with aviation/defence additions) --- it governs the organisation's
processes: document control, configuration management, traceability, supplier
control, corrective action. DO-178B is specifically about \emph{software}
development assurance. They're complementary: AS9100 is the quality framework the
company operates under (the SAP QM records, the controlled procedures), DO-178B
is the software-specific objective set within it. The resume's SAP QM /
AS9100 work is the QMS side; the DO-178B Level A work is the software-assurance
side.

\tierbanner{Tier 3 --- Trap \& depth (principal-engineer level)}%
{A paragraph plus the trap.}

\question{Q3.1}{You have 100\% MC/DC coverage and all requirements-based tests
pass. Can the software still be wrong? How?}
Yes --- coverage proves the \emph{code you wrote} was thoroughly exercised and
meets the \emph{requirements you wrote}; it says nothing about \emph{missing}
requirements or wrong requirements. If a requirement is absent (an unhandled edge
case nobody specified) or incorrect (it specifies the wrong behaviour), the code
can be fully covered, fully traceable, fully passing --- and still do the wrong
thing. That's the verification-vs-\emph{validation} gap. Also, coverage doesn't
catch concurrency/timing bugs (a race may need a specific interleaving), or
errors in the requirements-to-test interpretation. \trap the trap is treating
``100\% MC/DC + all tests green'' as ``proven correct.'' It proves the code is
\emph{right against the requirements}; validation, requirements review, robustness
(out-of-range/fault-injection) testing, and timing analysis cover what coverage
cannot. Naming that boundary is senior-level V\&V judgement.

\question{Q3.2}{Why does MC/DC use $N+1$ tests instead of all $2^N$
combinations, and what's the risk in the reduction?}
Exhaustive testing of $N$ conditions needs $2^N$ combinations --- infeasible for
large decisions. MC/DC achieves high confidence with ~$N+1$ tests by requiring,
\emph{per condition}, just the pair that demonstrates independent effect ---
which catches any single condition being wrong (stuck, inverted, or omitted)
without testing every combination. The risk in the reduction: MC/DC can miss
certain \emph{multiple-condition} interaction errors and is sensitive to how the
boolean expression is structured (coupled conditions, short-circuit evaluation),
and a poorly written compound condition can be hard to achieve MC/DC on
honestly. \trap the trap is thinking MC/DC = exhaustive; it's a carefully chosen
\emph{subset} that's sound for single-condition faults. The discipline is
keeping decisions simple (avoid deeply coupled conditions) so MC/DC is both
achievable and meaningful --- which also makes the code reviewable.

\question{Q3.3}{A defect is found in flight-critical code late in the program.
Walk the rigorous response, not just ``fix it.''}
Beyond the patch: (1) \textbf{contain} --- assess if shipped/other builds are
affected (configuration management tells you where the code is used); (2)
\textbf{root-cause} --- 5-why to the true origin, and ask whether the \emph{same
class} of defect exists elsewhere (a missing volatile, an unprotected shared
variable, a requirement gap); (3) \textbf{assess process escape} --- why did
verification not catch it? (missing test, missing requirement, inadequate
coverage, review gap) and fix the \emph{process}, not just the code; (4)
\textbf{fix with traceability} --- update requirement/design/code/test together,
re-run the specific and regression tests, re-measure coverage; (5)
\textbf{document} the corrective action (AS9100). \trap the trap is the
narrow ``fix the bug and move on''; in safety-critical work a defect is also
\emph{evidence of a process escape}, and the rigorous response asks ``what else
slipped through the same gap, and how do we close it?'' That systemic view ---
fixing the defect \emph{and} the hole that let it through --- is what distinguishes
real V\&V engineering, and it's the mindset behind the resume's defect-management
and traceability work.

\question{Q3.4}{How would you verify a requirement like ``the system shall fail
over to the backup channel within 50\,ms of a primary fault''?}
Decompose it into testable assertions and pick methods per assertion:
\textbf{requirements-based} HIL test --- inject a primary-channel fault (fault
injection) and \emph{measure} the failover time against the 50\,ms bound,
repeated across fault types and operating points, with the timing captured
objectively (a logged timestamp or a scope on a GPIO, Chapter 11); \textbf{robustness}
--- inject faults at boundary moments (mid-frame, during a CCDL exchange) to catch
timing races; \textbf{structural} --- confirm the failover code path is covered
(MC/DC on the detection logic) and that no path bypasses the deadline;
\textbf{analysis} --- a worst-case timing analysis (Chapter 9) showing the 50\,ms
is met even in the worst scheduling case, since a single test only samples one
interleaving. \trap the trap is a single happy-path test that passes once; a
\emph{timing} safety requirement needs measured failover under injected faults
\emph{plus} worst-case analysis, because the deadline must hold every time, not
on average. This is precisely the CCDL ``integrity, timing, failover''
verification on the resume, expressed as a rigorous V\&V argument.

%=====================================================================
\section{Coding Gym}
%=====================================================================

\begin{partnote}
Two host-compilable problems (verified under
\code{gcc -Wall -Wextra -std=c11}) for the computational side of V\&V: an MC/DC
independence-pair checker and a requirements-traceability coverage query --- the
logic behind the coverage and traceability tools.
\end{partnote}

%---------------------------------------------------------------
\begin{gymproblem}{MC/DC coverage checker}

\textbf{Problem.} Given a set of test vectors (input combinations and outcomes)
for a boolean decision, determine whether MC/DC is satisfied --- each condition
shown to independently affect the outcome.

\textbf{Hints.}
\begin{itemize}
\item For condition \code{c}, look for two tests differing only in bit \code{c}
  with different outcomes.
\item Full MC/DC requires such an independence pair for every condition.
\end{itemize}

\attempt

\textbf{Solution.}
\begin{cgym}
#include <stdint.h>
#include <stdio.h>

static int mcdc_condition_covered(const uint32_t *inmask, const int *outcome,
                                  int ntests, int cond) {
    for (int a = 0; a < ntests; a++)
        for (int b = 0; b < ntests; b++)
            if ((inmask[a] ^ inmask[b]) == (1u << cond) && outcome[a] != outcome[b])
                return 1;          /* independence pair found for this condition */
    return 0;
}
static int mcdc_full(const uint32_t *inmask, const int *outcome, int ntests, int nconds) {
    for (int c = 0; c < nconds; c++)
        if (!mcdc_condition_covered(inmask, outcome, ntests, c)) return 0;
    return 1;
}

int main(void) {
    /* decision = a && b ; bit0=a, bit1=b. Minimal MC/DC set: {11,10,01} */
    uint32_t good[3] = {0x3, 0x2, 0x1};  int gout[3] = {1, 0, 0};
    uint32_t weak[2] = {0x3, 0x0};       int wout[2] = {1, 0};   /* only 11,00 */
    printf("good=%d weak=%d\n",
           mcdc_full(good, gout, 3, 2), mcdc_full(weak, wout, 2, 2));   /* 1 0 */
    return 0;
}
\end{cgym}

The checker finds, for each condition, a pair of tests that differ in \emph{only}
that bit and yield different results --- the formal definition of MC/DC. The
``good'' 3-test set (\{11,10,01\}) achieves it for \code{a \&\& b}; the ``weak''
set (\{11,00\}) doesn't, because flipping both conditions at once can't show
either one \emph{independently}. This is the logic a coverage tool runs to certify
Level A, and writing it makes MC/DC concrete rather than a buzzword.

\followups
\begin{enumerate}
\item \emph{``Why does \{11,00\} fail MC/DC?''} $\rightarrow$ Both conditions
  change together; neither is shown to independently flip the outcome.
\item \emph{``How many tests for 3 conditions?''} $\rightarrow$ ~$N+1$ = 4 (vs 8
  exhaustive), if the conditions are independent.
\item \emph{``Short-circuit \code{\&\&} complicates this --- how?''}
  $\rightarrow$ Some combinations are unreachable, affecting which pairs are
  achievable.
\end{enumerate}
\end{gymproblem}

%---------------------------------------------------------------
\begin{gymproblem}{Requirements traceability coverage}

\textbf{Problem.} Given how many tests cover each requirement, compute the
coverage percentage and find the first untested requirement.

\textbf{Hints.}
\begin{itemize}
\item Covered = requirements with $\ge 1$ test; report the percentage and the
  first gap.
\end{itemize}

\attempt

\textbf{Solution.}
\begin{cgym}
#include <stdio.h>

static int trace_coverage_pct(const int *tests_per_req, int n) {
    int covered = 0;
    for (int i = 0; i < n; i++) if (tests_per_req[i] > 0) covered++;
    return covered * 100 / n;
}
static int first_uncovered(const int *tests_per_req, int n) {
    for (int i = 0; i < n; i++) if (tests_per_req[i] == 0) return i;
    return -1;       /* all covered */
}

int main(void) {
    int tpr[5] = {2, 1, 0, 3, 1};      /* requirement index 2 has no test */
    printf("coverage=%d%% first_gap=%d\n",
           trace_coverage_pct(tpr, 5), first_uncovered(tpr, 5));   /* 80% 2 */
    return 0;
}
\end{cgym}

This is the query at the heart of a traceability matrix: 4 of 5 requirements have
a test (80\%), and requirement 2 is the gap. In a certifiable process you drive
this to 100\% --- every requirement traced to at least one passing test (the
resume's claim) --- and the gap report tells you exactly what to write next. The
real tool (DOORS, or a managed SAP QM record) does this bidirectionally and links
to pass/fail results, but the logic is this simple coverage check.

\followups
\begin{enumerate}
\item \emph{``Also check the reverse direction.''} $\rightarrow$ Every \emph{test}
  must trace to a requirement --- flag orphan tests too.
\item \emph{``Weight by requirement criticality?''} $\rightarrow$ A Level-A
  requirement with no test is worse than a Level-D one --- prioritise gaps.
\item \emph{``Tie to pass/fail, not just existence.''} $\rightarrow$ A covered
  requirement whose test \emph{fails} isn't satisfied --- track results, not just
  links.
\end{enumerate}
\end{gymproblem}

%=====================================================================
\section{Link to Her Resume}
%=====================================================================

\begin{resumelink}
\textbf{DO-178B Level A V\&V:} ``I verified DFCC flight software against DO-178B
Level A requirements --- which means MC/DC structural coverage, requirements-based
testing, and 100\% requirements-to-test traceability. I can explain why Level A
needs MC/DC (branch coverage misses boolean-logic faults) and why coverage alone
doesn't prove correctness --- validation and robustness testing cover what it
can't.''

\medskip
\textbf{Test authoring \& traceability:} ``I author ATP/QTP procedures --- numbered
steps with expected results --- and maintain the traceability matrix so every
requirement maps to a passing test, managed in SAP QM under AS9100. A gap in the
matrix is a defect in the process, not just a missing test.''

\medskip
\textbf{Defect \& failover verification:} ``My defect work is root-cause-driven
(5-why to the true origin, then ask what else shares that cause) with regression
re-test. For a timing/failover requirement I verify with fault injection on HIL,
\emph{measure} the failover time, and back it with worst-case analysis --- not a
single happy-path pass.''

\medskip
\small Maps directly to the resume: DO-178B Level A verification, white/black-box
firmware testing, 100\% ATP/QTP traceability, SAP QM / AS9100, and CCDL
integrity/timing/failover verification. Deep project scripting is Chapter 15,
from \code{inputs/INTAKE.md} only.
\end{resumelink}

%=====================================================================
\section{Self-Test Scorecard}
%=====================================================================

\begin{scorecard}
\begin{enumerate}
\item Verification vs validation, in one line each.
\item Order statement / branch / MC/DC by strictness and say what each adds.
\item What does MC/DC require for each condition, and roughly how many tests?
\item What changes between DO-178B levels A, B, and C for coverage?
\item What is requirements traceability and what does 100\% mean?
\item ATP vs QTP?
\item What is HIL testing good for, and one limit?
\item Can fully-covered, all-passing software still be wrong? How?
\item Outline the defect lifecycle and where 5-why fits.
\item How would you verify a ``fail over within 50\,ms'' requirement rigorously?
\end{enumerate}
\end{scorecard}

\begin{answerkey}
\begin{enumerate}
\item Verification: built it right (meets requirements). Validation: built the
  right thing (meets the real need).
\item Statement (every line) $<$ branch (both outcomes of each decision) $<$
  MC/DC (each condition independently affects the outcome).
\item A test pair differing only in that condition, with different outcomes;
  ~$N+1$ tests for $N$ conditions.
\item MC/DC required at A, decision (branch) at B, statement at C --- stricter
  coverage for more severe failure consequences.
\item Bidirectional links requirement$\leftrightarrow$design$\leftrightarrow$
  code$\leftrightarrow$test; 100\% = every requirement has a passing test and
  every test traces to a requirement.
\item ATP accepts the product against requirements (numbered steps + expected
  results); QTP qualifies it against the full spec/environment.
\item Running real firmware against a real-time simulated plant for realistic/
  fault scenarios; limited by model fidelity (unmodelled effects escape).
\item Yes --- coverage proves it's right against the \emph{written} requirements;
  a missing/wrong requirement, a timing race, or a validation gap can still make
  it wrong.
\item Report $\rightarrow$ triage $\rightarrow$ root-cause (5-why) $\rightarrow$
  fix $\rightarrow$ re-test + regression $\rightarrow$ close, with traceability;
  5-why finds the true cause before fixing.
\item Fault-injection HIL test that \emph{measures} failover time across fault
  types/boundaries, plus MC/DC on the detection logic and a worst-case timing
  analysis --- not a single happy-path pass.
\end{enumerate}
\end{answerkey}

\begin{partnote}
\emph{End of Chapter 12, and of Part C (Systems \& Domain).} Next: Part D ---
the \textbf{Coding Gym}, 60 graded Embedded C interview problems, pure practice
for the coding rounds.
\end{partnote}
